Let \((\Omega,\mathscr F,P)\) carry a full law compatible with the maintained restrictions, and fix \(x\in\mathcal X\) with \(p_x>0\).  We first make the atomwise conditioning step explicit.  By the repaired declaration of \(\mathcal X\), its measurable structure is discrete.  Thus \(\{x\}\) is measurable, and measurability of the random variable \(X\) implies that
\[
A_x:=X^{-1}(\{x\})=\operatorname{xEvent}(x)\in\mathscr F.
\]
The observed carrier
\[
\mathcal O_{\mathrm{obs}}
=\mathcal X\times\{0,1\}^{3}\times(\mathcal Y\cup\{\star\})
\]
is also finite and discrete.  Write \(W=Y\) on \(\{S=1\}\) and \(W=\star\) on \(\{S=0\}\), so that
\[
\operatorname{observedDatum}(\omega)
=\bigl(X(\omega),Z(\omega),D(\omega),S(\omega),W(\omega)\bigr)
=O(\omega).
\]
For each singleton \(\{o\}\subseteq\mathcal O_{\mathrm{obs}}\), its inverse image under \(O\) is a finite intersection of measurable coordinate events.  In particular, the last-coordinate fibers are \(\{S=1,Y=i\}\) for \(i\in\mathcal Y\) and \(\{S=0\}\) for \(\star\).  Hence every singleton inverse image is measurable.  Every subset of the finite carrier is a finite union of singletons, so \(O=\operatorname{observedDatum}\) is measurable.  Consequently
\[
P_{\mathrm{obs}}=P\mathbin{\operatorname{map}}O,
\qquad
(P\mathbin{\operatorname{map}}O)(\mathcal O_{\mathrm{obs}})
=P\{O\in\mathcal O_{\mathrm{obs}}\}=1,
\]
and \(P_{\mathrm{obs}}\) is a probability law.  Moreover, for every measurable \(E\subseteq\mathcal O_{\mathrm{obs}}\), the defining identity for a measurable pushforward gives
\[
P_{\mathrm{obs}}(E)=P\bigl(O^{-1}(E)\bigr).
\tag{1}
\]
Taking \(E=\{o:X(o)=x\}\) in (1) shows \(P(A_x)=p_x\).

We next transport conditional instrument independence to the positive atom \(A_x\).  Put
\[
V=(D_0,D_1,S_0,S_1,Y_0,Y_1),
\qquad \mathscr G=\sigma(X).
\]
For any event \(A\in\sigma(V)\) and \(z\in\{0,1\}\), conditional instrument independence gives the conditional bundle identity
\[
E_P\!\left[\mathbf 1_A\mathbf 1_{\{Z=z\}}\mid\mathscr G\right]
=E_P\!\left[\mathbf 1_A\mid\mathscr G\right]
 E_P\!\left[\mathbf 1_{\{Z=z\}}\mid\mathscr G\right]
\quad P\text{-almost surely}.
\tag{2}
\]
Because \(X\) has a finite discrete carrier, on every atom \(A_u\) with \(P(A_u)>0\) the three conditional expectations in (2) have the ratio versions
\[
\frac{P(A\cap\{Z=z\}\cap A_u)}{P(A_u)},
\qquad
\frac{P(A\cap A_u)}{P(A_u)},
\qquad
\frac{P(\{Z=z\}\cap A_u)}{P(A_u)},
\]
respectively.  Evaluating (2) on the positive atom \(A_x\) therefore yields
\[
P(A\cap\{Z=z\}\cap A_x)P(A_x)
=P(A\cap A_x)P(\{Z=z\}\cap A_x).
\tag{3}
\]
Let \(\pi_1(x)=\pi(x)\) and \(\pi_0(x)=1-\pi(x)\).  Identity (1), the definition of \(\pi\), and instrument overlap imply
\[
P(\{Z=z\}\cap A_x)=p_x\pi_z(x)
\ge p_x\varepsilon_Z>0.
\tag{4}
\]
Dividing (3) by the two positive factors justified by \(p_x>0\) and (4) gives the required atomwise conditional-independence equality
\[
P(A\mid Z=z,X=x)
=\frac{P(A\cap\{Z=z\}\cap A_x)}{P(\{Z=z\}\cap A_x)}
=\frac{P(A\cap A_x)}{P(A_x)}
=P(A\mid X=x).
\tag{5}
\]

Fix \(i\in\mathcal Y\).  On \(\{Z=z,D=0,S=1\}\), treatment consistency gives \(D_z=0\), selection exclusion-consistency gives \(S_0=1\), and outcome exclusion-consistency gives \(Y_0=Y\).  Therefore (1) rewrites the relevant observed atom as
\[
\begin{aligned}
&P_{\mathrm{obs}}(Y=i,S=1,D=0\mid Z=z,X=x)\\
&\qquad=P(Y_0=i,S_0=1,D_z=0\mid Z=z,X=x)\\
&\qquad=P(Y_0=i,S_0=1,D_z=0\mid X=x),
\end{aligned}
\tag{6}
\]
where the last equality is (5) applied to the measurable event \(A=\{Y_0=i,S_0=1,D_z=0\}\in\sigma(V)\).  Substituting \(z=0\) and \(z=1\) into (6), then using the definition of \(\ell_i(x)\), gives
\[
\begin{aligned}
\ell_i(x)
&=P(Y_0=i,S_0=1,D_0=0\mid X=x)\\
&\quad-P(Y_0=i,S_0=1,D_1=0\mid X=x).
\end{aligned}
\tag{7}
\]
No defiers and binary treatment imply
\[
\{D_1=0\}\subseteq\{D_0=0\},
\qquad
\{D_0=0\}\setminus\{D_1=0\}=\{D_0=0,D_1=1\}=C.
\tag{8}
\]
Using (8) to subtract the nested events in (7) proves
\[
\ell_i(x)=P(Y_0=i,S_0=1,C\mid X=x)\ge0.
\tag{9}
\]

The higher-arm calculation is symmetric but its subtraction order is reversed.  On \(\{Z=z,D=1,S=1\}\), the three consistency and exclusion restrictions give \((D_z,S_1,Y_1)=(1,1,Y)\).  Thus (1) and (5) imply, for every \(j\in\mathcal Y\),
\[
P_{\mathrm{obs}}(Y=j,S=1,D=1\mid Z=z,X=x)
=P(Y_1=j,S_1=1,D_z=1\mid X=x).
\tag{10}
\]
Hence
\[
\begin{aligned}
h_j(x)
&=P(Y_1=j,S_1=1,D_1=1\mid X=x)\\
&\quad-P(Y_1=j,S_1=1,D_0=1\mid X=x).
\end{aligned}
\tag{11}
\]
Again by no defiers and binary treatment,
\[
\{D_0=1\}\subseteq\{D_1=1\},
\qquad
\{D_1=1\}\setminus\{D_0=1\}=C.
\tag{12}
\]
Subtracting the nested events in (11) therefore yields
\[
h_j(x)=P(Y_1=j,S_1=1,C\mid X=x)\ge0.
\tag{13}
\]

Because \(\mathcal Y=\{0,\ldots,K-1\}\) is finite and the events \(\{Y_a=k\}\), \(k\in\mathcal Y\), partition the relevant selected stratum, finite additivity in (9) and (13) gives
\[
q_0(x)=\sum_{i=0}^{K-1}\ell_i(x)=P(S_0=1,C\mid X=x),
\qquad
q_1(x)=\sum_{j=0}^{K-1}h_j(x)=P(S_1=1,C\mid X=x).
\tag{14}
\]
Subtracting the two equalities in (14) proves
\[
\Delta q(x)
=P(S_1=1,C\mid X=x)-P(S_0=1,C\mid X=x).
\tag{15}
\]

It remains to identify the survivor mass and read the direction.  If \(P(C\mid X=x)=0\), then (14) gives \(q_0(x)=q_1(x)=m(x)=0\), and also
\[
P(S_0=S_1=1,C\mid X=x)=0=m(x).
\tag{16}
\]
Suppose instead that \(P(C\mid X=x)>0\).  Under weak selection monotonicity, if \(d(x)=+1\), then \(S_1\ge S_0\) almost surely conditional on \((C,X=x)\).  Since selection is binary,
\[
\{S_0=1,C\}\subseteq\{S_1=1,C\}\quad\text{conditional on }X=x.
\]
Consequently, by (14),
\[
P(S_0=S_1=1,C\mid X=x)=q_0(x)=m(x)
\tag{17}
\]
and
\[
\Delta q(x)=P(S_0=0,S_1=1,C\mid X=x)\ge0.
\tag{18}
\]
If \(d(x)=-1\), weak selection monotonicity instead gives \(S_1\le S_0\) almost surely conditional on \((C,X=x)\), and the same nested-event argument gives
\[
P(S_0=S_1=1,C\mid X=x)=q_1(x)=m(x)
\tag{19}
\]
and
\[
\Delta q(x)=-P(S_0=1,S_1=0,C\mid X=x)\le0.
\tag{20}
\]
Equations (16), (17), and (19) establish the survivor-mass formula in every case.  Equation (20) rules out \(d(x)=-1\) when \(\Delta q(x)>0\), so a positive gap forces \(d(x)=+1\); equation (18) rules out \(d(x)=+1\) when \(\Delta q(x)<0\), so a negative gap forces \(d(x)=-1\).  Finally, when \(\Delta q(x)=0\), (18) or (20) makes the corresponding one-sided selection stratum null.  Thus \(S_0=S_1\) almost surely within \((C,X=x)\), and the same conditional latent law satisfies the weak-monotonicity inequality with either label \(d(x)=+1\) or \(d(x)=-1\).  Relabeling that direction leaves \(P_{\mathrm{obs}}\) unchanged; if the complier mass is zero, the restriction is vacuous and the same conclusion holds.  Hence the observable zero gap does not identify the direction.